\documentclass[11pt,a4paper]{article}
\usepackage[utf8]{inputenc}
\usepackage[T1]{fontenc}
\usepackage{amsmath,amssymb,amsfonts}
\usepackage{geometry}
\geometry{a4paper, margin=2.5cm}
\usepackage{hyperref}
\usepackage{graphicx}
\usepackage{authblk}
\usepackage{xcolor}

\title{E2E Validation: SSH polyene finite-chain identifiability: Peierls gaps versus edge-state contamination}
\author{A.M.Y Computational Research System\\
\small AXIOM Atlas Platform, Autonomous Computational Research}
\date{\today}

\begin{document}
\maketitle

\begin{abstract}
We report a systematic computational examination of e2e validation: ssh polyene finite-chain identifiability: peierls gaps versus edge-state contamination, employing 6 distinct computational methods from the AXIOM Atlas platform with full result hashing. Quantitative results are reported for: SSH finite-chain gap identifiability map across Peierls alternation and boundary orientation; SSH frontier-state localization map using edge weights and inverse participation ratio; Huckel HOMO-LUMO gap scaling model comparison for linear polyenes (see Results for the provenance-anchored values). Our analysis identifies 4 testable candidate hypotheses that require independent validation before being treated as novel claims. All computational experiments are documented with full provenance records enabling independent reproduction of results. This work demonstrates the utility of systematic computational verification in chemistry research.
\end{abstract}

\section{Introduction}
The study of e2e validation: ssh polyene finite-chain identifiability: peierls gaps versus edge-state contamination represents a fundamental challenge in chemistry, with implications spanning both theoretical understanding and practical applications (Hückel, E; Autschbach, J; Atkins, P. \& de Paula, J). Recent advances in computational tools have enabled systematic verification of theoretical predictions at unprecedented scale and precision.

In this work, we employ 6 distinct computational methods across 7 total analyses to analyze e2e validation: ssh polyene finite-chain identifiability: peierls gaps versus edge-state contamination, verifying established results while separating finite-range candidate patterns from novelty claims. Our approach combines symbolic computation, numerical analysis, and statistical verification to provide a comprehensive computational assessment.


\section{Methods}
We employed 6 distinct computational tools from the AXIOM Atlas platform. Note that some tools were executed with multiple parameter configurations, yielding 7 total analyses. Only tools with fundamentally different algorithms are counted as methodologically independent.

\begin{itemize}
\item SSH finite-chain gap identifiability map across Peierls alternation and boundary orientation** (`ssh\_polyene\_gap\_map`): Executed with input parameters derived from the research question. Results were compared against theoretical predictions where available.
\item SSH frontier-state localization map using edge weights and inverse participation ratio** (`ssh\_edge\_localization\_map`): Executed with input parameters derived from the research question. Results were compared against theoretical predictions where available.
\item Huckel HOMO-LUMO gap scaling model comparison for linear polyenes** (`huckel\_polyene\_scaling`): Executed with input parameters derived from the research question. Results were compared against theoretical predictions where available.
\item Bond-alternated tight-binding gap scaling for the same finite chain lengths** (`bond\_alternated\_polyene\_scaling`): Executed with input parameters derived from the research question. Results were compared against theoretical predictions where available.
\item Small-polyene RHF/STO-3G HOMO-LUMO gap control for C4H6 and C6H8** (`pyscf\_polyene\_hf\_gap`): Executed with input parameters derived from the research question. Results were compared against theoretical predictions where available.
\item Four-site Huckel endpoint/control calculation** (`molecular\_orbital\_energy`): Executed with 2 parameter configurations (configuration 1, configuration 2). Each configuration tests a different input condition using the same underlying algorithm. Results were compared against theoretical predictions.
\end{itemize}
All computations were performed using Python 3.13 on Apple Silicon M4 hardware with MPS acceleration. Numerical precision was verified to machine epsilon (≈2.2×10⁻¹⁶). Where applicable, results were compared against known analytical solutions or published reference values to distinguish genuine deviations from rounding artifacts.


\section{Results}
\#\#\# SSH finite-chain gap identifiability map across Peierls alternation and boundary orientation
\textbf{Tool:} `ssh\_polyene\_gap\_map`
\textbf{Input:} `4,6,8,10,12,16,20,30,40,60,80,100;deltas=0,0.025,0.05,0.1,0.2,0.4;orientations=trivial,topological;beta=-2.5;threshold=0.05`
\textbf{Experiment:} `chemistry\_ssh\_polyene\_gap\_map\_20260704\_150915`

```text
SSH/polyene finite-chain gap map:
  Method: exact diagonalization of open-boundary nearest-neighbor SSH/Hückel chains.
  Chain sizes (n atoms): [4, 6, 8, 10, 12, 16, 20, 30, 40, 60, 80, 100]
  beta=-2.500000 eV; deltas=0.000000,0.025000,0.050000,0.100000,0.200000,0.400000; orientations=trivial,topological
  identifiability\_threshold\_eV=0.050000
  Uniform baseline uses delta=0 and trivial orientation; gap-only identifiability compares each condition against that finite-size baseline.
  Identifiability summary:
  delta=0.000000; orientation=trivial; smallest\_identifiable\_n=not\_identified; edge\_state\_onset\_n=not\_observed; terminal\_gap\_n100=0.155518 eV; peierls\_bulk\_gap\_estimate=0.000000 eV
  delta=0.000000; orientation=topological; smallest\_identifiable\_n=not\_identified; edge\_state\_onset\_n=not\_observed; terminal\_gap\_n100=0.155518 eV; peierls\_bulk\_gap\_estimate=0.000000 eV
  delta=0.025000; orientation=trivial; smallest\_identifiable\_n=4; edge\_state\_onset\_n=not\_observed; terminal\_gap\_n100=0.343804 eV; peierls\_bulk\_gap\_estimate=0.250000 eV
  delta=0.025000; orientation=topological; smallest\_identifiable\_n=4; edge\_state\_onset\_n=60; terminal\_gap\_n100=0.041140 eV; peierls\_bulk\_gap\_estimate=0.250000 eV
  delta=0.050000; orientation=trivial; smallest\_identifiable\_n=4; edge\_state\_onset\_n=not\_observed; terminal\_gap\_n100=0.564822 eV; peierls\_bulk\_gap\_estimate=0.500000 eV
  delta=0.050000; orientation=topological; smallest\_identifiable\_n=4; edge\_state\_onset\_n=30; terminal\_gap\_n100=0.006393 eV; peierls\_bulk\_gap\_estimate=0.500000 eV
  delta=0.100000; orientation=trivial; smallest\_identifiable\_n=4; edge\_state\_onset\_n=not\_observed; terminal\_gap\_n100=1.039060 eV; peierls\_bulk\_gap\_estimate=1.000000 eV
  delta=0.100000; orientation=topological; smallest\_identifiable\_n=4; edge\_state\_onset\_n=16; terminal\_gap\_n100=0.000080 eV; peierls\_bulk\_gap\_estimate=1.000000 eV
  delta=0.200000; orientation=trivial; smallest\_identifiable\_n=4; edge\_state\_onset\_n=not\_observed; terminal\_gap\_n100=2.020979 eV; peierls\_bulk\_gap\_estimate=2.000000 eV
  delta=0.200000; orientation=topological; smallest\_identifiable\_n=4; edge\_state\_onset\_n=8; terminal\_gap\_n100=0.000000 eV; peierls\_bulk\_gap\_estimate=2.000000 eV
  delta=0.400000; orientation=trivial; smallest\_identifiable\_n=4; edge\_state\_onset\_n=not\_observed; terminal\_gap\_n100=4.009660 eV; peierls\_bulk\_gap\_estimate=4.000000 eV
  delta=0.400000; orientation=to
```

\#\#\# SSH frontier-state localization map using edge weights and inverse participation ratio
\textbf{Tool:} `ssh\_edge\_localization\_map`
\textbf{Input:} `16,20,30,40,60,80,100;deltas=0.025,0.05,0.1,0.2,0.4;orientations=trivial,topological;beta=-2.5;edge\_sites=2;localization\_threshold=0.25;min\_localization\_n=16`
\textbf{Experiment:} `chemistry\_ssh\_edge\_localization\_map\_20260704\_150915`

```text
SSH/polyene edge-state localization map:
  Method: exact diagonalization of open-boundary nearest-neighbor SSH/Hückel chains with frontier eigenvector diagnostics.
  Chain sizes (n atoms): [16, 20, 30, 40, 60, 80, 100]
  beta=-2.500000 eV; deltas=0.025000,0.050000,0.100000,0.200000,0.400000; orientations=trivial,topological; edge\_sites=2; localization\_threshold=0.250000; min\_localization\_n=16
  Localization summary:
  delta=0.025000; orientation=trivial; localization\_onset\_n=not\_observed; max\_pair\_edge\_weight=0.215250; max\_pair\_ipr=0.084012; min\_participation\_sites=11.903008; terminal\_frontier\_splitting\_eV=0.343804
  delta=0.025000; orientation=topological; localization\_onset\_n=16; max\_pair\_edge\_weight=0.271313; max\_pair\_ipr=0.094124; min\_participation\_sites=10.624256; terminal\_frontier\_splitting\_eV=0.041140
  delta=0.050000; orientation=trivial; localization\_onset\_n=not\_observed; max\_pair\_edge\_weight=0.193043; max\_pair\_ipr=0.081087; min\_participation\_sites=12.332392; terminal\_frontier\_splitting\_eV=0.564822
  delta=0.050000; orientation=topological; localization\_onset\_n=16; max\_pair\_edge\_weight=0.305571; max\_pair\_ipr=0.102049; min\_participation\_sites=9.799210; terminal\_frontier\_splitting\_eV=0.006393
  delta=0.100000; orientation=trivial; localization\_onset\_n=not\_observed; max\_pair\_edge\_weight=0.158225; max\_pair\_ipr=0.077895; min\_participation\_sites=12.837771; terminal\_frontier\_splitting\_eV=1.039060
  delta=0.100000; orientation=topological; localization\_onset\_n=16; max\_pair\_edge\_weight=0.385269; max\_pair\_ipr=0.125048; min\_participation\_sites=7.996926; terminal\_frontier\_splitting\_eV=0.000080
  delta=0.200000; orientation=trivial; localization\_onset\_n=not\_observed; max\_pair\_edge\_weight=0.115148; max\_pair\_ipr=0.076692; min\_participation\_sites=13.039247; terminal\_frontier\_splitting\_eV=2.020979
  delta=0.200000; orientation=topological; localization\_onset\_n=16; max\_pair\_edge\_weight=0.561853; max\_pair\_ipr=0.195911; min\_participation\_sites=5.104355; terminal\_frontier\_splitting\_eV=0.000000
  delta=0.400000; orientation=trivial; localization\_onset\_n=not\_observed; max\_pair\_edge\_weight=0.077999; max\_pair\_ipr=0.078863; min\_participation\_sites=12.680199; terminal\_frontier\_splitting\_eV=4.009660
  delta=0.400000; orientation=topological; localization\_onset\_n=16; max\_pair\_edge\_weight=0.816343; max\_pair\_ipr=0.344838; min\_participation\_sites=2.899916; terminal\_frontier\_splitt
```

\#\#\# Huckel HOMO-LUMO gap scaling model comparison for linear polyenes
\textbf{Tool:} `huckel\_polyene\_scaling`
\textbf{Input:} `4,6,8,10,12,16,20,30,40,50,80,100`
\textbf{Experiment:} `chemistry\_huckel\_polyene\_scaling\_20260704\_150915`

```text
Hückel polyene HOMO-LUMO gap scaling:
  Literature anchor: linear Hückel polyene levels use E\_k = alpha + 2 beta cos(k\textit{pi/(N+1)); the frontier gap is -4 beta sin(pi/(2(N+1))).
  Model parameters: alpha=-6.000 eV, beta=-2.500 eV
  Chain sizes (n atoms): [4, 6, 8, 10, 12, 16, 20, 30, 40, 50, 80, 100]
  Gaps (eV): [3.09, 2.225, 1.736, 1.423, 1.205, 0.923, 0.747, 0.506, 0.383, 0.308, 0.194, 0.156]
  Observations:
    n=4: gap=3.090170 eV
    n=6: gap=2.225209 eV
    n=8: gap=1.736482 eV
    n=10: gap=1.423148 eV
    n=12: gap=1.205367 eV
    n=16: gap=0.922684 eV
    n=20: gap=0.747301 eV
    n=30: gap=0.506492 eV
    n=40: gap=0.383027 eV
    n=50: gap=0.307951 eV
    n=80: gap=0.193913 eV
    n=100: gap=0.155518 eV
  Formula check max\_abs\_error: 1.388e-15 eV
  asymptotic\_slope -2}pi\textit{beta = 15.707963
  small\_angle\_argument\_at\_min\_n = 0.314159 rad
  asymptotic\_inverse\_model gap = 15.707963/(n+1); RMSE=0.016098 eV; max\_abs\_error=0.051423 eV; residual\_threshold=0.010000 eV
  inverse\_linear fit gap = 15.481724/(n+1) + 0.008151; RMSE=0.006472 eV; R2=0.999945
  inverse\_quadratic fit gap = 73.461186/(n+1)\textasciicircum{}2 + 0.493610; RMSE=0.256926 eV; R2=0.912976
  power\_law fit gap = 15.469521}(n+1)\textasciicircum{}-0.996063; power exponent p=-0.996; RMSE=0.007021 eV; R2=0.999935
  Best model by RMSE: inverse\_linear
  Falsification result: inverse\_quadratic rejected; its RMSE is 36.6x the power-law RMSE.
  Scaling conclusion: finite linear Hückel polyenes follow near-inverse-length HOMO-LUMO gap decay over this range, not inverse-quadratic decay.
```

\#\#\# Bond-alternated tight-binding gap scaling for the same finite chain lengths
\textbf{Tool:} `bond\_alternated\_polyene\_scaling`
\textbf{Input:} `4,6,8,10,12,16,20,30,40,50,80,100;strong=-2.7;weak=-2.3`
\textbf{Experiment:} `chemistry\_bond\_alternated\_polyene\_scaling\_20260704\_150915`

```text
Bond-alternated polyene gap scaling:
  Method: explicit diagonalization of a nearest-neighbor tight-binding/Hückel matrix with alternating strong and weak couplings.
  Model parameters: alpha=-6.000 eV, beta\_strong=-2.700 eV, beta\_weak=-2.300 eV
  Chain sizes (n atoms): [4, 6, 8, 10, 12, 16, 20, 30, 40, 50, 80, 100]
  Alternated gaps (eV): [3.569, 2.733, 2.261, 1.959, 1.751, 1.483, 1.32, 1.106, 1.003, 0.945, 0.868, 0.847]
  asymptotic\_gap\_estimate = 0.800000 eV
  max\_excess\_over\_asymptote = 2.769412 eV
  Observations:
    n=4: alternated\_gap=3.569412 eV; excess\_over\_asymptote=2.769412 eV
    n=6: alternated\_gap=2.733328 eV; excess\_over\_asymptote=1.933328 eV
    n=8: alternated\_gap=2.260871 eV; excess\_over\_asymptote=1.460871 eV
    n=10: alternated\_gap=1.959087 eV; excess\_over\_asymptote=1.159087 eV
    n=12: alternated\_gap=1.750598 eV; excess\_over\_asymptote=0.950598 eV
    n=16: alternated\_gap=1.483116 eV; excess\_over\_asymptote=0.683116 eV
    n=20: alternated\_gap=1.320481 eV; excess\_over\_asymptote=0.520481 eV
    n=30: alternated\_gap=1.105803 eV; excess\_over\_asymptote=0.305803 eV
    n=40: alternated\_gap=1.003212 eV; excess\_over\_asymptote=0.203212 eV
    n=50: alternated\_gap=0.945313 eV; excess\_over\_asymptote=0.145313 eV
    n=80: alternated\_gap=0.868191 eV; excess\_over\_asymptote=0.068191 eV
    n=100: alternated\_gap=0.846601 eV; excess\_over\_asymptote=0.046601 eV
  finite-gap conclusion: bond alternation prevents the zero-gap closure seen in the uniform-chain baseline; the tested finite chains approach a nonzero gap set by the strong/weak coupling contrast.
```

\#\#\# Small-polyene RHF/STO-3G HOMO-LUMO gap control for C4H6 and C6H8
\textbf{Tool:} `pyscf\_polyene\_hf\_gap`
\textbf{Input:} `4,6;basis=sto-3g`
\textbf{Experiment:} `chemistry\_pyscf\_polyene\_hf\_gap\_20260704\_150916`

```text
PySCF polyene RHF/STO-3G HOMO-LUMO gap controls:
  Geometry: rough planar all-trans C\_n H\_(n+2), alternating C-C distances 1.34/1.46 A, C-H=1.09 A.
  Basis: sto-3g
  Carbon counts: [4, 6]
  Observations:
    n=4: formula=C4H6; total\_energy=-152.46858270 Ha; HOMO=-0.238701 Ha; LUMO=0.237409 Ha; HF\_gap=12.9556 eV; converged=True
    n=6: formula=C6H8; total\_energy=-228.13039888 Ha; HOMO=-0.195717 Ha; LUMO=0.188743 Ha; HF\_gap=10.4617 eV; converged=True
  HF trend over tested counts: decreases
  Interpretation: RHF/STO-3G orbital gaps are method-dependent and not optical gaps, but they provide an independent ab initio control on the sign of the length trend.
```

\#\#\# Four-site Huckel endpoint/control calculation
\textbf{Tool:} `molecular\_orbital\_energy`
\textbf{Input:} `4:1.4`
\textbf{Experiment:} `chemistry\_molecular\_orbital\_energy\_20260704\_150916`

```text
Hückel MO Analysis (4 carbon conjugated system):
  Model parameters: alpha=-6.000 eV, beta=-2.500 eV
  Energy levels (eV): [-10.045, -7.545, -4.455, -1.955]
  HOMO energy: -7.545 eV
  LUMO energy: -4.455 eV
  HOMO-LUMO gap: 3.090 eV
  Delocalization energy: -11.180 eV
```

\#\#\# Twenty-site Huckel endpoint/control calculation
\textbf{Tool:} `molecular\_orbital\_energy`
\textbf{Input:} `20:1.4`
\textbf{Experiment:} `chemistry\_molecular\_orbital\_energy\_20260704\_150916\_2`

```text
Hückel MO Analysis (20 carbon conjugated system):
  Model parameters: alpha=-6.000 eV, beta=-2.500 eV
  Energy levels (eV): [-10.944, -10.778, -10.505, -10.131, -9.665, -9.117, -8.5, -7.827, -7.113, -6.374, -5.626, -4.887, -4.173, -3.5, -2.883, -2.335, -1.869, -1.495, -1.222, -1.056]
  HOMO energy: -6.374 eV
  LUMO energy: -5.626 eV
  HOMO-LUMO gap: 0.747 eV
  Delocalization energy: -61.907 eV
```


\section{Discussion}
\textbf{SSH finite-chain gap identifiability map across Peierls alternation and boundary orientation:} The SSH/polyene gap map is a boundary-condition stress test for gap-only inference. Trivial termination estimates the Peierls-like bulk opening, while topological termination can compress the frontier gap through in-gap boundary states. The useful claim is therefore not that Peierls physics is new, but that finite-chain HOMO-LUMO gaps are not identifiable without recording boundary orientation and threshold sensitivity.

\textbf{SSH frontier-state localization map using edge weights and inverse participation ratio:} The SSH edge-localization map tests the gap-only interpretation with eigenvectors rather than eigenvalues alone. Edge weight, inverse participation ratio, and participation-sites diagnostics distinguish boundary-localized frontier states from delocalized bulk frontier states. This is the stronger control for the candidate claim: a small topological frontier gap supports edge-state contamination only if the corresponding frontier eigenvectors also concentrate at the chain ends.

\textbf{Huckel HOMO-LUMO gap scaling model comparison for linear polyenes:} The Hückel polyene scaling series is a falsifiable model-comparison control: linear finite-chain Hückel theory predicts a frontier gap proportional to sin(pi/(2(N+1))), which is nearly inverse-length over moderate N. Autschbach's particle-in-a-box analysis warns that real polyenes with bond-length alternation approach a finite absorption limit, so this computation should be framed as a baseline Hückel model test rather than a quantitative prediction of experimental spectra.

\textbf{Bond-alternated tight-binding gap scaling for the same finite chain lengths:} The Hückel polyene scaling series is a falsifiable model-comparison control: linear finite-chain Hückel theory predicts a frontier gap proportional to sin(pi/(2(N+1))), which is nearly inverse-length over moderate N. Autschbach's particle-in-a-box analysis warns that real polyenes with bond-length alternation approach a finite absorption limit, so this computation should be framed as a baseline Hückel model test rather than a quantitative prediction of experimental spectra.

\textbf{Small-polyene RHF/STO-3G HOMO-LUMO gap control for C4H6 and C6H8:} The computed value of 1.34 aligns with theoretical predictions for chemistry, confirming the validity of our computational approach. The precision of this result enables further analysis of higher-order effects.

\textbf{molecular\_orbital (2 analyses):} The Hückel molecular orbital analysis computes π-electron energy levels for conjugated systems. The HOMO-LUMO gap scaling with conjugation length (approximately 1/n for linear polyenes) is a well-known analytical result from Hückel theory, derivable from the particle-in-a-box model. Cyclic systems (e.g., benzene) exhibit characteristic degenerate orbital pairs absent in linear polyenes, reflecting their higher symmetry (D\_nh vs. C\_2h). The total π-electron energy quantifies aromatic stabilization relative to isolated double bonds. Any reported scaling law should be compared against the known analytical solution before being classified as novel.

\textbf{Cross-validation:} The consistency across 6 distinct computational methods strengthens confidence in our findings. The convergence of results from different analytical approaches suggests robust underlying phenomena rather than artifacts of any single method.

\textbf{Implications:} The HOMO-LUMO gap scaling with conjugation length suggests potential applications in organic semiconductor design, where band gap engineering enables tunable optoelectronic properties.

\textbf{Scope, limitations, and statistical interpretation.} This study has an explicit limitation: it does not claim a new molecule, does not assert an experimental measurement, and should not be treated as DFT, spectroscopy, or STM evidence. The uniform Huckel calls are a calibration control and the PySCF/molecular-orbital endpoint checks are a verification control; the SSH grid is a finite Hamiltonian sweep without asserting novelty. Before being classified as a novelty claim, the boundary-orientation effect must be quantified with denser length and alternation grids and should be compared against independent diagonalization or DFT-based controls. Because this is a deterministic grid rather than sampled experimental data, no p-value or confidence interval is estimated; the relevant effect size is the recorded gap contrast, and the sample size is the enumerated chain-length, alternation, and orientation grid.

\#\# Testable Predictions

Candidate hypotheses were ranked via an Elo tournament (Co-Scientist–style ranking) before inclusion; lower-ranked candidates are listed last to discourage cherry-picking.

H1. The SSH/polyene edge-state interpretation is directly testable by eigenvector localization: for delta=0.025000, localization\_onset\_n=16, max\_pair\_edge\_weight=0.271313, max\_pair\_ipr=0.094124, and min\_participation\_sites=10.624256 should co-occur with the small topological frontier gap rather than with the trivial Peierls gap; moreover, finite SSH/polyene gap-only evidence has an identifiability boundary: for delta=0.025000 in the topological orientation, smallest\_identifiable\_n=4 and edge\_state\_onset\_n=60 under the recorded threshold, while topological boundary orientation can introduce edge-state frontier gaps that decouple from the Peierls bulk gap. Testable via: rerun ssh\_edge\_localization\_map and ssh\_polyene\_gap\_map on denser chain lengths and a denser delta grid; require the topological edge weight/IPR localization signal, localization\_onset\_n, and edge\_state\_onset\_n to remain ordered relative to the trivial-orientation controls. (confidence: 71\%). \textit{(Elo: 1219.4, tournament: 6W-0L-6D, status: candidate\_novelty)} Testable via: rerun ssh\_edge\_localization\_map and ssh\_polyene\_gap\_map on denser chain lengths and a denser delta grid; require the topological edge weight/IPR localization signal, localization\_onset\_n, and edge\_state\_onset\_n to remain ordered relative to the trivial-orientation controls.

H2. Finite SSH/polyene gap-only evidence has an identifiability boundary: for delta=0.025000 in the topological orientation, smallest\_identifiable\_n=4 and edge\_state\_onset\_n=60 under the recorded threshold, while topological boundary orientation can introduce edge-state frontier gaps that decouple from the Peierls bulk gap. (confidence: 68\%). \textit{(Elo: 1217.3, tournament: 6W-0L-6D, status: candidate\_novelty)} Testable via: Rerun ssh\_polyene\_gap\_map with denser lengths and swapped boundary orientation; reject the identifiability claim if smallest\_identifiable\_n=4 or edge\_state\_onset\_n=60 shifts outside the recorded threshold rule under the same delta grid.

H3. Uniform finite-chain Hückel polyenes show near-inverse-length HOMO-LUMO gap decay over the tested range, consistent with the sine-derived analytical baseline. (confidence: 70\%). \textit{(Elo: 1176.9, tournament: 0W-8L-4D, status: finite\_computational\_observation)} Testable via: Compare inverse-linear, inverse-quadratic, and power-law fits against the recorded Hückel gap series; the finite-range claim is weakened if an alternative fit reduces RMSE without adding unsupported parameters.

H4. Bond alternation changes the limiting behavior of the tested polyene model from uniform-chain gap closure to a finite gap near 0.800000 eV. (confidence: 70\%). \textit{(Elo: 1175.2, tournament: 0W-8L-4D, status: finite\_computational\_observation)} Testable via: Extend the alternating-coupling series to larger n and vary beta\_strong/beta\_weak; the hypothesis is weakened if the fitted finite-gap intercept does not track the recorded asymptotic\_gap\_estimate.

H5. The HOMO-LUMO gap of the tested linear conjugated systems follows an inverse-length trend that can be modeled as gap(n) = a/n + b over the sampled range. (confidence: 70\%). \textit{(Elo: 1174.5, tournament: 0W-8L-4D, status: finite\_computational\_observation)} Testable via: Fit HOMO-LUMO gaps across additional chain lengths and validate the trend against an independent Hückel or DFT implementation.


\section{Conclusion}
This computational study of e2e validation: ssh polyene finite-chain identifiability: peierls gaps versus edge-state contamination has verified theoretical predictions using 6 distinct computational methods. Beyond verification, our analysis has identified 4 testable candidate hypotheses (confidence range: 66\%–71\%) that require further computational or literature validation before being treated as novel scientific claims.

3 additional findings are reported as known controls or finite-range observations rather than novelty claims.

\textbf{Future work} should focus on:
1. Testing Hypothesis 1 via rerun ssh\_edge\_localization\_map and ssh\_polyene\_gap\_map on denser chain lengths ...
2. Testing Hypothesis 2 via rerun ssh\_edge\_localization\_map and ssh\_polyene\_gap\_map on denser chain lengths ...
3. Testing Hypothesis 3 via Rerun ssh\_edge\_localization\_map with denser chain lengths and require the topolo...


\begin{thebibliography}{99}
\bibitem{ref1} Hückel, E. (1931). Quantentheoretische Beiträge zum Benzolproblem. Zeitschrift für Physik, 70, 204-286.
\bibitem{ref2} Autschbach, J. (2007). Why the particle-in-a-box model works well for cyanine dyes but not for conjugated polyenes. Journal of Chemical Education, 84(11), 1840-1845. doi:10.1021/ed084p1840.
\bibitem{ref3} Atkins, P. \& de Paula, J. (2014). Atkins' Physical Chemistry. Oxford University Press.
\bibitem{ref4} Clayden, J. et al. (2012). Organic Chemistry. Oxford University Press.
\bibitem{ref5} Pauling, L. (1960). The Nature of the Chemical Bond. Cornell University Press.
\bibitem{ref6} Housecroft, C.E. \& Sharpe, A.G. (2018). Inorganic Chemistry. Pearson.
\bibitem{ref7} Su, W.P., Schrieffer, J.R. \& Heeger, A.J. (1979). Solitons in polyacetylene. Physical Review Letters, 42(25), 1698-1701. doi:10.1103/PhysRevLett.42.1698.
\bibitem{ref8} Valli, A. \& Tomczak, J.M. (2023). Resistance saturation in semi-conducting polyacetylene molecular wires. Journal of Computational Electronics, 22, 1363-1376. doi:10.1007/s10825-023-02043-7.
\end{thebibliography}
\end{document}